\documentclass[11pt]{article}

%####### DON'T CHANGE MARGIN SETTINGS ###########
\newcommand{\keywordfont}{\textsc}
\newcommand{\keyword}[1]{%
  \marginpar{\raggedright\small\keywordfont{#1}}}
\reversemarginpar
\usepackage[a4paper, top=2.5cm, bottom=2.5cm, outer=2cm, inner=3.3cm, marginparwidth=72pt, heightrounded]{geometry}
%#################################

\usepackage{amsmath}
\usepackage{amssymb}
\usepackage{amsthm}
\usepackage{hyperref}
\usepackage{graphicx}
\usepackage{microtype}
\usepackage{float}

\newtheorem{claim}{Claim}
\newtheorem{conjecture}{Conjecture}




\begin{document}

\Large
\begin{center}
\textbf{MA3K7 Assignment $3$}
\\
Alexandros Sergiou (5505826)
\end{center}
\normalsize

\section{Entry}

\keyword{I KNOW}From the assumption in the problem,
\begin{itemize}
\item Two players are given a pair of positive integers $(n,s)$ corresponding to $n$ heaps and $s$ stones
\item Players take turns placing one stone in one of $n$ heaps (empty before the game starts) 
\item During the Nim phase, players take turns removing any positive number of stones from a single heap. The player who removes the final stone loses.
\end{itemize}
I assume that Player 1 makes the first placement and that turns continue alternating between the two phases. Therefore, if $s$ is even, Player 1
begins the Nim phase, whereas if $s$ is odd, Player 2 begins.

\medskip
\keyword{NOTATION}
For brevity, I define:

\begin{itemize}
    
\item $H=(h_1,h_2,\dots,h_n)$ corresponds to the resulting heaps, where $h_i$ is the number of stones in heap $i$
\item The Nim sum of the heaps, defined by
$h_1 \oplus h_2 \oplus \cdots \oplus h_n,$
where $\oplus$ denotes binary addition without carrying.
\end{itemize}
\medskip
\keyword{I WANT} I want to investigate this game and how it's played while exploring it with python, specifically I will:
\begin{itemize}
\item Simulate the game to understand how it works.    
\item Explore how different values of $(n,s)$ affect the winning strategy.
\item Investigate the usefulness of the Nim sum.
\item Explore how changing the rules affects the outcome.
\end{itemize}

\section{Attack}
To begin, I used Python to simulate the game to help me visualise it. I followed a similar approach to the one seen in the lectures. 


\medskip
\keyword{OBSERVE}
The code helped me understand how the game works. However, since both players made random moves, the results did not tell me which player had a winning strategy.

\medskip
\keyword{IDEA}
I will now try to improve the simulation by allowing both players to make optimal moves, rather than choosing them randomly. This should help me investigate how the values
of $n$ and $s$ affect the outcome.

\medskip
\keyword{CONJECTURE}
From my initial computations, I wondered
whether the winning player depended on
whether $s$ was odd or even, or on the
relationship between $n$ and $s$.
To investigate this, I tested several
different values of $(n,s)$ using Python.



\pagebreak

\medskip
\keyword{CHECK} 
\begin{table}[h]
\centering
\begin{tabular}{|c|c|c|}
\hline
  heaps ($n$) &  stones ($s$) & Winner \\
\hline
2 & 3 & Player 2 \\
\hline
2 & 4 & Player 2 \\
\hline
5 & 4 & Player 1 \\
\hline
5 & 3 & Player 2 \\
\hline
6 & 5 & Player 2 \\
\hline
6 & 6 & Player 2 \\
\hline
7 & 6 & Player 2 \\
\hline
7 & 1 & Player 1 \\
\hline
\end{tabular}
\caption{Winning players for different values of $n$ and $s$,
assuming optimal play.}
\end{table}



\keyword{STUCK}
After comparing the results of the eight games, I could not identify a pattern based on the parity of $n$ and $s$, or whether $n>s$. Although the Python code identified which player could force a win, the results did not explain the strategies behind these outcomes.

\medskip
\keyword{IDEA}
Rather than testing more values of $(n,s)$,
I decided to investigate whether the Nim sum could help explain the winning strategies.

\medskip
\keyword{ATTEMPT}
I first considered the case $n=2$ and $s=4$. Since there are only two heaps, I wondered whether Player 2 could respond to each of Player 1's moves by placing a stone in the opposite heap.
Suppose Player 1 places the first stone in
heap 1. Player 2 can respond by placing a
stone in heap 2, giving $(1,1)$.
Whatever heap Player 1 chooses next, Player 2 can again place a stone in the opposite heap. This guarantees the final configuration $(2,2)$.
Now since $2\oplus2=0$, and Player 1 begins the second phase, Player 2 can force a win under the last-stone-loses rule.

\medskip
\keyword{CHECK}
I then considered whether the same strategy would work for other values of $s$.
For $n=2$ and $s=6$, Player 2 can again place a stone in the opposite heap after every move made by Player 1. This guarantees the final
configuration $(3,3)$, for which
$$ 3\oplus3=0. $$
Since Player 1 begins the Nim phase, Player 2 can force a win.



\medskip
\keyword{OBSERVE}
When $s$ is odd, Player 1 makes the final
placement, so Player 2 cannot necessarily
leave the heaps equal. However, Player 2
then begins the Nim phase, which suggests
that a different strategy may be possible.


\keyword{CHECK}

To investigate this, I considered the case $n=2$ and $s=5$.
Using the same strategy as before, Player 2 can ensure that after four placements the heaps are $(2,2)$. Player 1 then places the final stone, giving either $(3,2)$ or $(2,3)$. In both cases, the Nim sum is
$$ 3\oplus2=1\neq0.$$
Since Player 2 begins the Nim phase, this is a winning position for Player 2.

\smallskip
\keyword{NOTE}
This suggests that the important factor is not simply whether the Nim sum is zero, but also which player begins the second phase.
I then considered the case $n=2$ and $s=2$,
to see whether the same strategy would work
when fewer stones were available.
Player 2 can place their stone in the
opposite heap, producing $(1,1)$, whose
Nim sum is
$$1\oplus1=0.$$
However, unlike the cases $s=4$ and $s=6$,
this is a winning position for Player 1
under the last-stone-loses rule.
Player 1 can remove one stone, leaving
Player 2 forced to take the final stone
and lose.
This made me realise that the usual
Nim-sum criterion does not always apply
under the last-stone-loses rule.

\keyword{AHA}

It appears we can use the nim sum to determine winning and losing positions
under the last-stone-loses rule, provided
at least one heap contains more than one stone.
However, when every heap contains at most
one stone and at least one stone remains,
the winning criterion is reversed.
In this case, a Nim sum of zero indicates
a winning position for the player whose
turn it is.
I shall check this observation
using Python.



\keyword{CHECK}

To verify my observations, I wrote another
Python function that examines all possible
moves without using the Nim sum.
I used this to determine winning and
losing positions in the second phase,
and compared the results with those
obtained using my original
Nim-sum-based algorithm.

\begin{table}[h]
\centering
\begin{tabular}{|c|c|c|c|}
\hline
Heaps & Nim sum & Recursive result & Outputs agree \\
\hline
$(2,2)$ & $0$ & Losing & Yes \\
$(1,1)$ & $0$ & Winning & Yes \\
$(3,2)$ & $1$ & Winning & Yes \\
$(1,1,1)$ & $1$ & Losing & Yes \\
$(1,1,1,1)$ & $0$ & Winning & Yes \\
$(3,3)$ & $0$ & Losing & Yes \\
\hline
\end{tabular}
\caption{Comparison of recursive results with the
mis\`ere Nim winning criterion.}
\end{table}


\keyword{NOTE}
The results agreed for all six configurations, supporting my earlier observation about the
Nim sum and the exception under the
last-stone-loses rule.


\section{Solution}
From my Python experiments, I found that although there was no obvious pattern for every pair $(n,s)$, I could identify some cases where a player has a winning strategy.
I will now use my observations about the Nim sum to justify these strategies.


\begin{claim}
If $s>n$ and $s$ is odd, then Player 2 has a winning strategy
\end{claim}

\begin{proof}
Since $s>n$, there are more stones than heaps, so at least one heap must contain more than one stone after the first phase.
Furthermore, since $s$ is odd, there must
be an odd number of odd-sized heaps.
As the least significant binary digit of
the Nim sum is determined by the parity
of the odd-sized heaps, the Nim sum
must be non-zero.
Therefore, the final configuration has
at least one heap containing more than
one stone and a non-zero Nim sum.
Under the last-stone-loses rule, this
is a winning position for the player
whose turn it is.
Since $s$ is odd, Player 2 begins the
second phase and can therefore force a win.
\end{proof}


\keyword{CHECK}
I then considered whether the condition $s>n$ was necessary for Player 2 to win.
For example, when $(n,s)=(5,3)$, it is possible for the final configuration to be $(1,1,1,0,0)$. Although its Nim sum is non-zero, this is a losing position for Player 2 under the last-stone-loses rule.

However, Player 2 can prevent this configuration. After Player 1's first placement, Player 2 can place a stone in the same heap, ensuring that at least one heap contains more than one stone.

Regardless of Player 1's final placement, the Nim sum must be non-zero, since the total number of stones is odd.
As Player 2 begins the second phase, they
can therefore force a win.

\keyword{OBSERVE}

This shows that the condition $s>n$ is
sufficient, but not necessary, for
Player 2 to have a winning strategy
when $s$ is odd.
Although the values of $(n,s)$ can guarantee
a favourable Nim position in some cases,
in others the players must use their
placements strategically to achieve one.

\begin{claim}
Suppose at least one stone remains and every heap contains at most one stone. Under the last-stone-loses rule, a position is winning for the player whose turn it is if and only if its Nim sum is zero.
\end{claim}

\begin{proof}
Since every non-empty heap contains exactly one stone, a player can only remove one stone per turn.
Let $k$ be the number of remaining stones.
If $k$ is even, the second player will take
the final stone and lose, so the player
moving first has a winning position.
If $k$ is odd, the first player will take
the final stone and lose, so the position
is losing for the player whose turn it is.
Since each heap contains either $0$ or $1$
stone, the Nim sum is

$$ h_1\oplus\cdots\oplus h_n = \begin{cases}
0, & \text{if } k \text{ is even},\\
1, & \text{if } k \text{ is odd}. \end{cases}$$
Therefore, the position is winning for
the player whose turn it is if and only
if its Nim sum is zero.
\end{proof}



\section{Review}

\keyword{CHECK}
To check my findings, I used Python to test
some additional values of $(n,s)$ satisfying
the conditions of my first claim.
For each case, the program identified Player 2 as the winner, agreeing with my mathematical argument.

I also tested configurations where every heap contained at most one stone. The results agreed with my second claim, including configurations containing empty heaps.



\keyword{REFLECT}
After initially running my Python simulations, I thought I had identified a reliable pattern. However, after testing more cases, I realised that my initial observations were not sufficient to explain the results. This led me to investigate the usefulness of the Nim sum, where I discovered that it does not behave in the same way under the last-stone-loses rule as in standard Nim. I found this interesting, as it led me to the result established in my second claim in the Solution section. As the objective of this assignment was to explore the game rather than find a general solution, I experimented with various scenarios and strategies to see whether I could identify any patterns. In the following section, I will explore how variations in the rules affect the winning strategies.


\keyword{EXTENSION}
I then decided to investigate how changing
the last-stone-loses rule would affect the
winning strategies. In particular, I considered
the variation where the player who removes
the final stone wins, essentially turning
the second phase into standard Nim while
keeping the first phase unchanged.


\pagebreak

\keyword{NOTICE}
I modified my Python code to compare the
winning players under both variations,
assuming optimal play. Some of the
results are shown below.

\begin{table}[h]
\centering
\begin{tabular}{|c|c|c|c|}
\hline
$n$ & $s$ & Mis\`ere Nim & Standard Nim \\
\hline
2 & 1 & Player 1 & Player 2 \\
2 & 2 & Player 1 & Player 2 \\
2 & 3 & Player 2 & Player 2 \\
2 & 4 & Player 2 & Player 2 \\
3 & 1 & Player 1 & Player 2 \\
3 & 2 & Player 1 & Player 2 \\
3 & 3 & Player 2 & Player 2 \\
3 & 4 & Player 1 & Player 1 \\
4 & 2 & Player 1 & Player 2 \\
4 & 3 & Player 2 & Player 2 \\
4 & 4 & Player 1 & Player 2 \\
4 & 5 & Player 2 & Player 2 \\

\hline
\end{tabular}
\caption{Comparison of the winning players
under the two different rules.}
\end{table}

\keyword{OBSERVE}
A couple of interesting observations can be made after running a python code comparing the winner before and after the rule change. The first: It appears that if $s>n$ the winning player was the same under both
variations.This can be explained using the fact that if $s>n$, at least one heap
must contain more than one stone after the
first phase. In that case, standard Nim and
mis\`ere Nim have the same winning criterion, so the outcome of the complete two-phase game is unchanged.I also noticed that whenever $s$ was odd and $s\geq 3$, Player 2 won under both rules in
all the cases I tested. 


\keyword{IDEA}
I wondered whether this would hold for
any number of heaps.
Suppose $s\geq 3$ is odd. After Player 1
places the first stone, Player 2 can place
their stone in the same heap.
This guarantees that at least one heap
contains more than one stone after the
first phase, regardless of the remaining
placements. Since $s$ is odd, the final Nim sum must be non-zero. Furthermore, Player 2 begins the second phase.
Therefore, Player 2 can force a win under
both standard Nim and the last-stone-loses
variation, for any $n\geq 1$.

\medskip
\keyword{EXTENSION}
I then decided to investigate another variation of the game, this time restricting the number of stones that could be removed during the second phase.
Instead of allowing players to remove any
number of stones from a single heap, I
considered the case where they could remove
at most two stones per turn.
I kept the last-stone-loses rule and the
first phase unchanged, to see how this
restriction would affect the winning strategies.

\medskip
\keyword{OBSERVE}
From the results, I noticed that restricting
the number of stones that could be removed
changed the winning player in several cases.
For example, when $(n,s)=(3,4)$, Player 1
wins under the original rules, whereas
Player 2 wins when players can remove
at most two stones per turn.
I decided to investigate why this happens.

\smallskip
\keyword{JUSTIFY}
Suppose Player 1 places the first stone in
heap 1. Player 2 can respond by placing a
stone in the same heap, producing $(2,0,0)$.
If Player 1 places the third stone in the
same heap, Player 2 can also place their
stone there, giving $(4,0,0)$.
Alternatively, if Player 1 places the third
stone in an empty heap, Player 2 can add
another stone to the first heap, producing
$(3,1,0)$, up to the order of the heaps.

I then checked why these two configurations
were losing positions.
From $(4,0,0)$, Player 1 can only leave
$(3,0,0)$ or $(2,0,0)$. In either case,
Player 2 can remove enough stones to
leave Player 1 with the final stone.
Similarly, from $(3,1,0)$, Player 1
can only leave $(2,1,0)$, $(1,1,0)$
or $(3,0,0)$.
In each case, Player 2 can make a move
that leaves exactly one stone remaining,
forcing Player 1 to take the last
stone and lose.
Therefore, both $(4,0,0)$ and
$(3,1,0)$ are losing positions for
Player 1, confirming the results
obtained using Python.
Hence, Player 2 can force a win
when $(n,s)=(3,4)$ under the
restricted rules.


\smallskip
\keyword{REFLECT}
What I found particularly interesting was
that restricting the possible moves changed
the usefulness of the Nim sum.
For example, $(4,0,0)$ has a non-zero Nim
sum and is winning under the original
rules. However, with a maximum removal of
two stones, it becomes a losing position.
This shows that the winning strategies
depend not only on the final heap
configuration, but also on which moves
are permitted during the second phase.





\end{document}
